\chapter{跨工况迁移诊断的基准脉络}

\section{任务定义与统一任务边界}
本文研究的核心任务是跨工况无监督领域自适应故障诊断。设源域训练样本为
\begin{equation}
\mathcal{D}_s=\{(x_i^s,y_i^s)\}_{i=1}^{n_s},
\end{equation}
其中 $x_i^s$ 为源域样本，$y_i^s\in\mathcal{Y}$ 为故障类别标签；目标域训练样本为
\begin{equation}
\mathcal{D}_t=\{x_j^t\}_{j=1}^{n_t},
\end{equation}
其中目标域标签在训练阶段不可用。本文的目标是在利用源域有标签样本和目标域无标签样本的条件下，学习一个诊断模型
\begin{equation}
f(x)=h(\phi(x)),
\end{equation}
使其在目标域测试集上的风险尽可能小，其中 $\phi(\cdot)$ 表示特征提取器，$h(\cdot)$ 表示分类头。

若将目标域真实风险记为
\begin{equation}
R_t(f)=\mathbb{E}_{(x,y)\sim P_t(x,y)}[\ell(f(x),y)],
\end{equation}
则无监督领域自适应的困难在于训练阶段无法直接最小化 $R_t(f)$。因此，实际训练通常采用
\begin{equation}
\mathcal{L}=\mathcal{L}_{cls}^{s}+\lambda\mathcal{L}_{adapt}+\mathcal{L}_{aux}
\end{equation}
的形式，其中 $\mathcal{L}_{cls}^{s}$ 为源域监督分类损失，$\mathcal{L}_{adapt}$ 为领域对齐损失，$\mathcal{L}_{aux}$ 为伪标签、一致性、原型或其他辅助约束项。

虽然本文的方法推导主要采用单源到单目标记号，但当前 benchmark 同时包含少量多源场景。对多源到单目标任务，可将源域集合记为
\begin{equation}
\mathcal{D}_{s}^{multi}=\{\mathcal{D}_{s}^{(m)}\}_{m=1}^{M},
\end{equation}
其中 $M\ge 2$ 表示源域个数。训练时可将多个源域小批量分别采样后在统一骨干上进行融合，从而在不改变主要记号体系的前提下兼容多源场景。

\section{基准构建原则}
与只追求某个模型单次最优结果不同，本文首先把 benchmark 视为研究工作的基础设施。一个可信的跨工况迁移诊断 benchmark，至少应满足以下四项原则。

\begin{enumerate}
\item \textbf{任务边界明确。} 必须清楚区分源域、目标域、是否允许使用目标域标签、是否固定折划分、是否包含多源场景等设置，避免不同论文之间“名义相同、协议不同”的比较误差。
\item \textbf{实现条件统一。} 不同方法应尽量共享相同的数据输入形式、归一化方式、骨干网络和结果导出结构，使比较聚焦于自适应机制，而非数据处理或骨干差异。
\item \textbf{模型选择可解释。} 无监督领域自适应缺少目标域验证标签，checkpoint 选择与早停策略本身就会影响结论，因此必须将其纳入实验协议而非隐含处理。
\item \textbf{结果产物可追踪。} 每次运行应统一保存指标、分析中间量、图件和 review 摘要，便于后续批量汇总、论文回填和误差追查。
\end{enumerate}

本文的 benchmark 设计正是围绕上述原则展开：在共享 FCN 骨干与统一数据协议下，将 11 个迁移场景和 8 种方法纳入同一 fixed-fold 实验计划，并通过统一脚手架输出 \texttt{result.json}、\texttt{review.json}、混淆矩阵、t-SNE 图以及汇总表格。

\section{域定义与统一输入表示}
在本文实验设定中，不同 TEP 操作模式被视为不同领域，即
\begin{equation}
\mathcal{M}=\{\mathrm{mode1},\mathrm{mode2},\ldots,\mathrm{mode6}\}.
\end{equation}
源域与目标域共享相同的标签空间，
\begin{equation}
\mathcal{Y}=\{0,1,\ldots,28\},
\end{equation}
但样本分布并不一致。输入样本采用多变量过程时序片段表示，原始形状为 $(T,C)$，其中 $T=600$ 表示时间窗口长度，$C=34$ 表示过程变量通道数；进入网络前转置为 channels-first 形式：
\begin{equation}
x^\prime=\mathrm{Transpose}(x)\in\mathbb{R}^{C\times T}.
\end{equation}
模型输出为 29 类工况与故障类别的预测概率。

当前 88-run benchmark 采用“两簇场景”的组织思路：一簇为 mode3 与 mode6 的互迁，另一簇为 mode1、mode2、mode5 之间的互迁与多源组合。这样的设计并不试图穷尽全部 30 个有向单源场景，而是在算力与实验广度之间取得折中，使 benchmark 既能覆盖不同迁移难度，也能在当前论文周期内完成运行与分析。

\section{共享骨干网络}
为降低骨干结构差异对实验结论的干扰，本文所有方法共享同一套一维全卷积网络（FCN）编码器和统一分类头。编码器由三层一维卷积构成，其计算过程可写为
\begin{align}
z^{(1)} &= \mathrm{ReLU}(\mathrm{Norm}(\mathrm{Conv1D}_{9,128}(x))),\\
z^{(2)} &= \mathrm{ReLU}(\mathrm{Norm}(\mathrm{Conv1D}_{5,256}(z^{(1)}))),\\
z^{(3)} &= \mathrm{ReLU}(\mathrm{Norm}(\mathrm{Conv1D}_{3,128}(z^{(2)}))),
\end{align}
其中卷积核长度依次为 9、5 和 3，输出通道数依次为 128、256 和 128，当前实现默认采用实例归一化。随后对时间维执行全局平均池化，得到紧凑特征向量
\begin{equation}
f=\frac{1}{T^\prime}\sum_{t=1}^{T^\prime} z_{:,t}^{(3)}\in\mathbb{R}^{128}.
\end{equation}
分类头采用一层隐藏维度为 128 的全连接网络：
\begin{equation}
p(x)=\mathrm{softmax}(W_2\,\mathrm{ReLU}(W_1 f+b_1)+b_2).
\end{equation}

该结构具备三个优点。其一，一维卷积对多变量时间窗口的局部动态模式具有较强建模能力；其二，参数规模适中，便于在 88 次 benchmark 运行中控制计算成本；其三，结构简单、变量少，更适合作为公平比较的公共底座。

\section{基础训练目标}
源域监督分类损失统一采用交叉熵形式：
\begin{equation}
\mathcal{L}_{cls}^{s}=
-\frac{1}{n_s}\sum_{i=1}^{n_s}\sum_{k=1}^{K}
\mathbf{1}[y_i^s=k]\log p_k(x_i^s),
\end{equation}
其中 $K=29$，$p_k(x)$ 为模型对第 $k$ 类的预测概率。不同领域自适应方法在此基础上分别引入协方差对齐、MMD 对齐、域对抗损失、条件对抗损失或最优传输损失；本文提出的 RCTA 则进一步加入伪标签、一致性和原型约束。

在训练脚手架层面，本文统一采用 Adam 优化器，基础学习率保持在 $10^{-4}$ 数量级，权重衰减设为 0；对于对抗式或更敏感的方法，额外采用梯度裁剪以降低训练波动。需要指出的是，本文不再把“所有方法完全同轮数训练”作为僵硬前提，而是允许在统一协议下通过 experiment config 对不同方法设置更合理的训练轮数和权重，这种做法更符合 benchmark 的实际目标，即比较方法机制而非人为限制超参数。

\section{典型基线方法与参考方法}
本节对当前 benchmark 涉及的典型方法进行简要说明，具体实验结果统一在第七章呈现。

\subsection{Source Only}
Source Only 仅依赖源域监督分类损失训练模型，不显式处理源域与目标域之间的分布偏移，其优化目标可写为
\begin{equation}
\mathcal{L}_{SO}=\mathcal{L}_{cls}^{s}.
\end{equation}
该方法构成所有领域自适应方法的基础对照组。若某一领域自适应方法在目标域上不能稳定优于 Source Only，则说明其对齐策略尚未真正缓解跨工况迁移带来的分布差异问题。

\subsection{Target Only}
Target Only 直接利用目标域标注样本进行监督训练，其目标为
\begin{equation}
\mathcal{L}_{TO}=\mathcal{L}_{cls}^{t}.
\end{equation}
该方法不属于无监督领域自适应范畴，而是作为\textbf{监督参考上界}纳入 benchmark，用于帮助判断跨域方法距离“目标域可监督训练”还有多大差距。将其与真正的无监督方法混为同类 baseline 并不合适，因此后文会明确标注其参考性质。

\subsection{分布对齐方法}
CORAL 通过约束源域与目标域特征协方差矩阵的一致性减小二阶统计差异~\cite{sun2016deep}。设源域与目标域特征分别为 $F_s\in\mathbb{R}^{n_s\times d}$ 和 $F_t\in\mathbb{R}^{n_t\times d}$，对应协方差矩阵为 $C_s,C_t$，则 CORAL 损失为
\begin{equation}
\mathcal{L}_{coral}=
\frac{\|C_s-C_t\|_F^2}{4d^2}.
\end{equation}

DAN 使用多核最大均值差异约束源域与目标域特征分布，在核空间中实现边缘分布对齐~\cite{long2015dan}。给定核映射 $\varphi(\cdot)$，其损失可表示为
\begin{equation}
\mathcal{L}_{mmd}=
\left\|
\frac{1}{n_s}\sum_{i=1}^{n_s}\varphi(f_i^s)-
\frac{1}{n_t}\sum_{j=1}^{n_t}\varphi(f_j^t)
\right\|_{\mathcal{H}}^{2}.
\end{equation}
这两类方法结构较直接、训练相对稳定，因此适合作为全局统计对齐的参考。

\subsection{对抗式方法}
DANN 在共享特征提取网络之后加入域判别器，并利用梯度反转层实现特征表示与域分类之间的对抗学习，从而逼近域不变特征~\cite{ganin2016dann}。设域标签 $d=1$ 表示源域，$d=0$ 表示目标域，域判别器为 $D(\cdot)$，则对抗损失可写为
\begin{equation}
\mathcal{L}_{adv}=
-\frac{1}{n_s}\sum_{i=1}^{n_s}\log D(f_i^s)
-\frac{1}{n_t}\sum_{j=1}^{n_t}\log(1-D(f_j^t)).
\end{equation}

CDAN 在 DANN 基础上进一步引入类别条件信息，使域判别器感知特征与类别预测的联合结构~\cite{long2018conditional}。若记分类器输出概率为 $p(x)$，则可构造条件映射
\begin{equation}
g(x)=f(x)\otimes p(x),
\end{equation}
其中 $\otimes$ 表示多线性组合。对抗式方法通常具有较强适应能力，但训练稳定性和超参数敏感性也更值得关注。

\subsection{最优传输方法}
DeepJDOT 通过联合建模特征距离与类别预测差异，构建源域与目标域样本之间的运输代价，在深度特征学习过程中同时优化分类与跨域匹配~\cite{damodaran2018deepjdot}。设 $\gamma$ 为源域和目标域样本间的运输矩阵，则其代价函数可写为
\begin{equation}
\mathcal{L}_{OT}=
\min_{\gamma\in\Pi(a,b)}
\sum_{i,j}\gamma_{ij}
\left(
\alpha\|f_i^s-f_j^t\|_2^2+
\beta\ell(y_i^s,p_j^t)
\right).
\end{equation}
该方法能够提供样本级跨域关系约束，但计算开销、批大小敏感性和代价尺度设置都更值得注意。

\section{评价指标与模型选择}
本文重点比较各方法的目标域诊断性能。评价指标主要包括目标域准确率、目标域平衡准确率、混淆矩阵，以及反映域对齐程度和类别结构清晰度的特征可视化结果。目标域准确率定义为
\begin{equation}
\mathrm{Acc}_{t}=\frac{1}{N_t^{eval}}\sum_{i=1}^{N_t^{eval}}\mathbf{1}[\hat{y}_i=y_i],
\end{equation}
平衡准确率定义为
\begin{equation}
\mathrm{BA}_{t}=\frac{1}{K}\sum_{k=1}^{K}
\frac{\mathrm{TP}_k}{\mathrm{TP}_k+\mathrm{FN}_k},
\end{equation}
其中 $\mathrm{TP}_k$ 和 $\mathrm{FN}_k$ 分别表示第 $k$ 类的真正例和假负例数量。相较于普通准确率，平衡准确率更能稳定反映多类别任务下不同类别的平均识别水平。

由于无监督领域自适应缺少可直接使用的目标域验证标签，checkpoint 选择与早停本身就是 benchmark 协议的重要组成部分。当前 benchmark 中，多数方法采用“源域评估准确率 + 目标域平均置信度”的混合选择指标：
\begin{equation}
S_{\mathrm{hyb}}=
w_s\cdot \mathrm{Acc}_{s}^{eval}
+w_c\cdot \bar{q}_t,
\end{equation}
其中 $\bar{q}_t$ 表示目标域训练集的平均预测置信度。对于更容易出现目标域过度自信的 RCTA，还会进一步加入过度自信惩罚与域判别偏离惩罚，这一点将在第六章展开说明。将模型选择规则显式写入 experiment config，而不是在训练后凭经验挑选结果，是本文 benchmark 设计中的关键环节。

\section{本章小结}
本章围绕后续方法与实验所需的共同前提，说明了跨工况迁移诊断的任务定义、benchmark 构建原则、域定义与输入表示、共享骨干网络、基础训练目标、典型基线方法和评价指标。与将“相关基础”孤立成章的写法不同，本章更强调后续方法设计与实验比较所依赖的统一前提，从而为第三章到第七章建立共同语境。
